\documentclass[12pt,a4paper,onecolumn]{article}
\title{Digest of The New York Times}
\author{The New York Times}
% \definecolor{dark-green}{rgb}{0,0.25,0}
% \usepackage[utf8]{inputenc}

\usepackage{fontspec,xunicode,xltxtra}
\setmainfont[Ligatures=Required,Mapping=tex-text]{Georgia}
% \setmainfont[Ligatures=Required,Mapping=tex-text]{Palatino Linotype}
\setsansfont[Ligatures=Required,Mapping=tex-text]{Myriad Pro}
\setmonofont[Ligatures=Required,Mapping=tex-text]{Courier New}
\usepackage{xeCJK}
\setCJKmainfont{宋体}
\setCJKmonofont{黑体}

% \usepackage{indentfirst}


\usepackage[dvips,dvipsnames,svgnames]{xcolor}
\definecolor{light-gray}{gray}{0.95}

%生成PDF的链接
\usepackage[dvipdfmx]{hyperref}
\hypersetup{
    bookmarks=true,         % show bookmarks bar?
    bookmarksopen=true,
    pdfpagemode=UseNone,    % options: UseNode, UseThumbs, UseOutlines, FullScreen
    pdfstartview=FitH,
    pdfborder=1,
    pdfhighlight=/P,
    pdfauthor={wuxch},
    unicode=true,           % non-Latin characters in Acrobat’s bookmarks
    colorlinks,             % false: boxed links; true: colored links
    linkcolor=blue,         % color of internal links
    citecolor=blue,        % color of links to bibliography
    filecolor=magenta,      % color of file links
    urlcolor=cyan           % color of external links
}
% 使用目录
\usepackage{makeidx}
% 使用图形
\usepackage{graphicx}


\usepackage{titletoc}
\dottedcontents{section}[1.5em]{}{2.3em}{1pc}

% % 页眉页脚的设定
% \usepackage{fancyhdr}
% \pagestyle{fancy}
% \lhead{\includegraphics[height=0.2in]{zte_logo}}

% \usepackage{footmisc}
\usepackage{pfnote}
\usepackage{fnpos}

% % 嵌入的代码显示
% \usepackage{listings}
% \lstset{language=C++, breaklines, extendedchars=false}
% \lstset{keywordstyle=\color{blue},
%         commentstyle=\color{SeaGreen},
%         stringstyle=\ttfamily,
%         showstringspaces=false,
%         backgroundcolor=\color{light-gray}}
% \lstset{basicstyle=\small}

\usepackage{booktabs}

\usepackage{tipa}
\usepackage{phonetic}
\usepackage[gen]{eurosym}

% \usepackage{CJK}
\usepackage{makeidx}
% \usepackage{glossary}

\usepackage{fancyhdr}

\usepackage{lastpage}
\cfoot{\textsf{\thepage\ of \pageref*{LastPage}}}

\pagestyle{fancy}
\fancyhead{}
\renewcommand{\headrulewidth}{0pt}
\pagenumbering{arabic}

\usepackage{graphicx}
\usepackage{picinpar}

\setlength{\parindent}{0in}
\setlength{\parskip}{2.5ex plus4mm minus3mm}
% \setlength{\parskip}{.3\baselineskip}

\setcounter{section}{0}

% bibtex
\usepackage[super]{natbib}
% \usepackage[globalcitecopy,labelstoglobalaux]{bibunits}
% \makeatletter
% \renewcommand{\@bibunitname}{\jobname.\the\@bibunitauxcnt}
% \makeatother

% \renewcommand{\refname}{缩略语}
% \renewcommand{\refname}{}

\makeatletter                   % Reference list option change
\renewcommand\@biblabel[1]{#1:  } % from [1] to 1: 
\makeatother                    %

\newcommand{\citenumfont}[1]{\textsf{#1}}

% \usepackage{tipa}
% \usepackage{phonetic}

\newcommand{\phonetic}[1]{{\fontspec[Colour=CC00CC]{Microsoft Sans Serif} #1}}


%%% Local Variables:
%%% TeX-master: "medical_trivia"
%%% End:

\makeindex
\begin{document}
\date{}
\thispagestyle{empty}
\renewcommand\contentsname{\textsf{Digest of The New York Times}}
\textsf{\tableofcontents}
\pagebreak
\setcounter{page}{1}
\section{China's Commodity Buying Spree}

HONG KONG -- Strong buying by China has helped lift commodity prices around the world this spring,
but growing evidence suggests that a sizable portion of this buying has been to build stockpiles in
China, and may not be sustainable.

At least 90 large freighters full of iron ore are idling off Chinese ports, where they face waits of
up to two weeks to unload because port storage operations are overflowing, chief executives of
shipping companies said in interviews this week. Yet actual steel production from that iron ore is
recovering much more slowly in China, and Chinese steel exports remain weak.

Commodities and shipping executives describe Chinese stockpiling in recent months of a range of
other commodities as well, including \\ aluminum, copper, nickel, tin, zinc, canola and soybeans.
Starting in April, China began stockpiling significant quantities of crude oil.

China's goals vary by commodity. Chinese companies have bought iron ore heavily on the spot market
in anticipation of higher prices in annual contract talks now nearing completion. The Chinese
government has been stockpiling oil and some metals for strategic reasons, and bought huge
quantities of aluminum and canola to insulate domestic producers of these goods from falling global
prices over the winter.

``There has been enormous stockpiling of all commodities'' by China, and this cannot continue
indefinitely, said Tim Huxley, the chief executive of Wah Kwong Maritime Transport Holdings, a big
shipping line based here.

Those extra purchases beyond China's daily needs have helped reverse the price collapse in
commodities that followed the economic downturn, but could also limit the scale of the rebound.

Moody's Investors Service announced on Wednesday that it was putting a negative outlook on the base
metals, mining and steel industries in Asia and the Pacific, having previously done so for these
sectors elsewhere.

``China's strategic stockpiling and replacement of lower-quality domestic production with
higher-quality imports have supported the recent rally in prices for many base metals, but we will
not see a sustainable turnaround in demand until the major economies of the U.S., Europe, and Japan
recover,'' said Terry Fanous, a senior vice president in Sydney for Moody's, adding that the leading
economies were not likely to recover until next year.

In the latest sign of weak overseas demand, the Chinese government announced on Thursday morning
that the country's exports fell 26.4 percent in May from a year earlier. Imports were down 25.2
percent from a fairly weak level a year ago, as China's overall trade surplus continued to narrow,
to \$13.39 billion.

The Standard \& Poor's GSCI, an index of global commodity prices, has risen 42 percent from its low
on Feb.~18, but is still less than half its record, set on July 3.

One of the best leading indicators of international trade in commodities is the Baltic Exchange Dry
Index, which measures the daily cost of chartering a large freighter. While the GSCI has continued
to rise in the last week, the freight index has fallen by a fifth in that period.

Richard S.~Elman, the chief executive of the Noble Group, Asia's largest diversified commodities
trading company, bounced up from the conference table in his office here when asked about freight
rates during an interview on Tuesday morning. He walked over to his desk, dominated by three
computer screens that partly obscure a perfect view of Hong Kong's harbor, and quickly punched up on
one screen a list of daily charter rates for large bulk carrier freighters.

The list showed ship owners charging \$58,000 a day now but just \\\$24,000 a day for charters next
year or in 2011 -- an indication that there will be more ships than cargoes in the years ahead,
particularly with shipyards still finishing vessels ordered during the recent boom.

Pointing to the rates for the next two years, he said, ``That's the real market'' for ships.

From an immense new sugar mill in Brazil to an extensive coking coal operation in Australia, Noble
is active in commodities around the globe, and its stock has nearly quadrupled since its low on
Oct.~24. Mr.~Elman voiced optimism about the future of the Chinese economy and of worldwide demand
for commodities, but cautioned that for some commodities, ``the futures prices have gone ahead'' of
the prices for physical delivery.

According to J.P. Morgan, China's iron ore imports were 33 percent higher in April than a year
earlier. Crude oil imports were up nearly 14 percent, aluminum oxide imports climbed 16 percent and
refined copper imports jumped 148 percent.

Imports of coal soared 168 percent as Chinese utilities bought more foreign coal while trying to
negotiate better prices with domestic producers.

Determining the percentage of each commodity being stockpiled is difficult, especially in China,
where scant data are released. Assessing steel demand, in particular, has become a subject of almost
obsessive interest among many shipping executives and economists as a barometer of emerging markets'
health and as an indicator of demand for things like iron ore and cars.

Sanjay Mehta, one of the four managing directors of Essar Global, the big Indian multinational in
steel, shipping and other heavy industries, estimated that North American steel mills were operating
at 50 to 60 percent of capacity, Chinese steel mills at 70 percent of capacity and Indian steel
mills at 100 percent of capacity.

The resilience of the Indian economy is helping to sustain demand for commodities, he said. But he
was cautious about the global economy. He suggested that part of China's purchasing over the last
several months represented an effort to rebuild inventories that were drawn down during the autumn
and winter.

``It is not all related to consumption,'' he said, predicting that prices would stay roughly at
current levels through the middle of 2011. Prices of many commodities have jumped sharply in recent
months -- spot oil prices, in particular, have doubled since late December. That is driving up the
price of gasoline and diesel in many countries.

Steel demand in China is already recovering for types of steel used in construction, Mr.~Elman said.
Local, provincial and national government agencies are ramping up investments quickly as part of
economic stimulus programs.

But demand has been slower to rebound for higher grades of steel used in consumer products, despite
\$1 billion in Chinese government incentives for the purchase of cars and household appliances.

Some economists say they are bullish on commodities because they believe that the United States and
European economies are on their way to recovery.

``The commodity price rally is for real,'' said Ajay Kapur, the chief global strategist at Mirae
Asset, a big Korean financial firm. ``I'm not expecting any huge correction from here.''

Other executives, particularly in shipping, are less optimistic, and see signs of a bubble in
freight rates, and possibly commodities, that may repeat the sudden rise and fall of prices last
year.

``The past two weeks have been nuts and, rather than cheering this sudden comeback of the dry bulk
market, I do have a considerable amount of concern that we are seeing the same bubble again,''
Kenneth Koo, the chairman and chief executive of the Tai Chong Cheang Steamship Company, another big
Hong Kong shipping line, wrote in an e-mail message. ``And like that past bubble, it's not going to
sustain.''

\section{China Faces Criticism Over New Software Censor}

BEIJING -- China is facing a storm of protest at home and abroad over new regulations requiring all
personal computers sold in the country to include software that can filter out pornography and other
``vulgar'' content from the Internet.

Industry executives, free-speech advocates and many computer users have reacted angrily to the new
mandate, which gives manufacturers until July 1 to preinstall the software on millions of new
machines.

``The compulsory installation of filtering software is a whole lot of useless flopping about,'' said
an editorial in The Wuhan Evening News.

Computer makers in the United States say it will be impossible to fulfill the requirement by the end
of the month and have asked the Chinese government to reconsider the directive. They say it raises
thorny questions about censorship and whether manufacturers will be liable if the software --
designed by a company with ties to China's military and public security agencies -- conflicts with
operating systems or causes computers to crash.

``To be honest, nobody really knows what this software is capable of,'' said one executive at an
American computer maker who spoke on the condition of anonymity because his company was still trying
to work with the government on the issue. Computer experts fear that it could allow the government
to monitor Internet use and collect personal information.

So far, the government has not shown any signs of backing down. On Tuesday, state-owned media and
officials defended the new software, \\known as Green Dam, and said suggestions that it could be
turned into government-operated spyware were exaggerated.

``If you have children or are expecting a child, you could understand the concerns of parents over
unhealthy online content,'' Qin Gang, a Foreign Ministry spokesman, said at a regular news
conference Tuesday.

Bloggers, dissidents and even some normally cautious state media outlets have greeted the
announcement with skepticism, questioning the software's reliability and wondering whether it could
be used to censor nonsexual content. Some criticized the government's decision to spend about \$6
million on a program that was not solicited through an open bidding process.

China vigorously restricts Internet content, regularly blocking Web \\pages devoted to Tibetan
independence, human rights issues and other politically sensitive subjects. An anti-vulgarity
campaign this year has shut down thousands of pornographic sites.

Green Dam's designers say the program, which uses image recognition technology and text filtering to
block material, is not capable of acting as spyware. Most important, they say the program can be
disabled or erased.

The Beijing News devoted a full page to Green Dam and found its performance less than stellar. For
example, an innocent math question that included the word ``balls'' was filtered out.

Writing on Green Dam's own Web site, one teacher said sexually explicit images slipped through and
another complained that it would not let her view pictures of pigs.

In an editorial he wrote for The Oriental Morning Post, Wang Lin, an associate professor at Hainan
University Law School, said the government should have consulted computer users and allowed other
companies to submit comparable programs. ``They've made a decision affecting tens of millions of
people without regard to their opinions,'' he said in an interview. ``It's like you buying a car and
the government telling you where you can drive.''

Executives at computer makers, which last year sold about 40 million PCs in China, agree with such
sentiments, although none were willing to speak on the record for fear of angering the government as
they press their case behind the scenes to have the rule reconsidered.

The government presented the regulations, drawn up without industry input, to PC makers on May 19,
but they were not released publicly until Monday.

Preliminary tests by software technicians at a number of companies have revealed significant
software tics that could affect a computer's operating system or other programs, according to the
American computer executive who spoke anonymously.

On Tuesday, a coalition of United States trade associations issued a statement that gently made
their case.

``We believe there should be an open and healthy dialogue on how \\parental control software can be
offered in the market in ways that ensure privacy, system reliability, freedom of expression, the
free flow of information, security and user choice,'' read the statement, which was signed by groups
including the Software \& Information Industry Association and the Information Technology Industry
Council, whose members include \\Lenovo, Dell, Apple and Hewlett-Packard.

While bloggers have been among the most incensed by the new rules, one, Wang Xiaofeng, said he was
not overly concerned the software would have a lasting impact on Internet freedom.

``People will always find a way to break through these controls,'' he said. ``It's just a shame the
government has to keep spending taxpayer money to build them.''

\section{China Protests Moving of Detainees}

BEIJING -- The Chinese government protested Thursday over the decision by the American government to
resettle a group of Chinese Muslims to the isolated archipelago of Palau. The Uighur men, former
detainees at Guant\'anamo Bay, have been in a state of limbo since they were cleared of wrongdoing
last fall.

In a news conference in Beijing, a Foreign Ministry spokesman, Qin Gang, called the detainees
``terrorist suspects'' and demanded they be sent back to China as soon as possible. He said that the
United States was ignoring international laws on terrorism in not doing so.

Shortly afterward, the United States announced that it had sent four Chinese Muslims who had also
been detained at Guant\'anamo to the island of Bermuda.

But hours after the detainees arrived on the island on Thursday, there was already the first sign
that their peaceful resettlement would not be without hitches. Bermuda is a British overseas
territory, but its premier, Ewart F.~Brown, did not mention to the British government that was
planning to accept the detainees.

The result was a terse British statement, as well as a pledge that the Foreign Office would help
Bermuda ``carry out a security assessment of the men.''

``We have underlined to the Bermuda government that it should have consulted the UK on whether this
falls within their competence, or is a foreign affairs or security issue for which the Bermuda
Government do not have delegated responsibility,'' a Foreign Office spokesperson said.

The conflict involves 17 Chinese Muslims -- of the Uighur ethnic group from northwest China -- who
were seized in Afghanistan during the American-led invasion in 2001 and accused of being terrorists,
but who maintained that they had been seeking refuge from Chinese persecution. China maintains that
they are members of an Islamic separatist group that is seeking independence for Xinjiang, a vast
province in northwest China whose Muslim population has often bridled under Chinese rule.

During the Bush administration, officials determined that the Uighurs did not travel to Afghanistan
with the intent to take any hostile action against the United States, and in the end, cleared them
for release.

But the United States resisted moving them onto American soil, and at one point, the Justice
Department argued they should never be admitted into the country because they ``sought to wage
terror'' in China.

In February, a federal appeals court overturned a district judge's order that would have freed the
men in the United States, saying the judge was assuming powers reserved to the president and
Congress.

That decision left in place the Bush administration's concession that the men were not enemy
combatants, the classification the government used to detain men at Guant\'anamo. The result was a
legal limbo.

American officials resisted sending the men back to China, fearing they would be jailed, tortured or
executed, and have been searching for other countries willing to take them. State Department
officials had said some 100 countries refused to take the men.

The resettlement of the four Uighurs to Bermuda was announced Thursday on the Justice Department's
Web site. It is the first time since 2006 that the government has resettled any of the Guant\'anamo
Uighurs. The Justice Department said it had conducted a comprehensive review of the four, including
a threat evaluation, before approving them for resettlement.

In 2006, five were transferred to Albania; there have been ``no reports of post-resettlement
engagement in criminal behavior or terrorist activities,'' the Justice Department said in a
statement.

``By helping accomplish the president's objective of closing Guant\'anamo, the transfer of these
detainees will make America safer,'' Attorney General Eric Holder said in the statement, adding that
the department was ``extremely grateful to the government of Bermuda for its assistance.''

Bermuda's premier, Ewart F.~Brown, said in a statement on Thursday that the four detainees had been
granted short-term asylum in Bermuda as a humanitarian gesture and a token of the island's
friendship with the United States. Washington, he said, will bear the costs associated with
relocating the men on the island.

``Of the prisoners held there,'' Mr.~Brown said, referring to Guant\'anamo, ``many are innocent men,
held without trial or any form of due process; many are refugees from their own lands whose
political views are contrary to the regimes in power. They have committed no crime,'' he said.

The detainees ``will have the opportunity to become naturalized citizens,'' he said, ``and
thereafter afforded the right to travel and leave Bermuda, potentially settling elsewhere.''

But, because of its colonial status, Mr.~Brown also said that it is ``important for everyone to
understand that this process is not complete.''

Mr.~Brown said that he had met with the British-appointed governor of Bermuda Thursday morning to
seek formal approval of the government's plan, and that the governor ``is seeking to further assess
the ramifications of this move before allowing the Government of Bermuda to fully implement this
action.''

The lawyers of one of the freed detainees, Abdul Nasser, said in a statement upon the men's arrival
on the sun-swept island that the men were grateful for Bermuda's action, the Agence France-Presse
reported.

``Growing up under communism, we always dreamed of living in peace and working in free society like
this one,'' Mr.~Nasser said. ``Today you have let freedom ring.''

Regarding the detainees now in Bermuda, Britain said it would ``decide on further steps as
appropriate,'' and stressed that the four ``do not have travel documents and will therefore not be
able to travel to the UK.''

China did not specify what, if any actions it would take to protest the relocation of the Uighurs.
Palau, a Pacific island nation and former United States territory, does not have diplomatic
relations with China, which may help insulate it from retribution.

\section{China's Economy Takes a Sharp Domestic Turn}

HONG KONG -- Since China entered the World Trade Organization in November 2001, a rising tide of
exports, combined with a torrent of investment, has lifted the country's economy ever higher, while
consumer spending has lagged.

But now, the Chinese economy relies increasingly on growth at home, as data released Thursday made
clear. A decline in exports has become a serious drag on growth, while government spending has led
domestic investments higher at a remarkable pace and consumer spending appears to have been fairly
strong as well.

Some economists wonder whether China is actually becoming too reliant on investment spending and
whether the government's economic stimulus program may be making this worse.

``For China's nascent economic recovery to be sustainable beyond the short term, policy makers must
take steps to ensure that consumption remains on a firm growth trajectory and that the investment
boom does not exacerbate the economy's structural imbalances,'' the chairwoman of China equities at
JPMorgan, Jing Ulrich, said in a research note.

Chinese exports plunged by a record 26.4 percent in May from a year earlier, the Chinese customs
agency announced Thursday, as buyers in industrialized countries remained cautious about placing
orders.

But investments in fixed assets like roads, factories and apartment buildings set a record in the
opposite direction.

Chinese investment expenditures rose 32.9 percent in the first five months of this year, compared
with the investments in the period last year, the National Bureau of Statistics announced in
Beijing.

Yu Song, a Goldman Sachs economist, calculated that after adjusting for inflation, Chinese
investment spending had grown in May at a breakneck pace, rising close to 50 percent from May of
last year.

The government's stimulus program is powering much of that increase, with spending on railroads
soaring 110.9 percent in the first five months of 2009, compared with the same period last year.

The Chinese media reported that retail sales in May, scheduled to be announced Friday morning, are
likely to show an increase of 15.2 percent. That would represent a modest acceleration from April,
when the increase from a year earlier was 14.8 percent, and a robust gain when adjusted for the
gradual decline in overall prices in China.

The Chinese government has struggled to keep economic data secret, and leaks to the media have
frequently been right.

One of the biggest supports for retail sales in China has been the auto market, with car sales
rising briskly as the government has offered various subsidies, especially in rural areas.

The big question for China is how long the economy can stay strong without the support of vigorous
exports. Wages and profits have slipped in export-oriented coastal zones of China, limiting the
spending power of consumers.

Chinese imports fell a little more slowly than exports last month, dropping 25.2 percent from the
period a year earlier. The Chinese trade surplus last month was \$13.39 billion, compared with
\$13.14 billion in April and \$16.37 billion a year ago.

\section{U.S. Takes On the Insular G.M. Culture}

DETROIT -- Fiat will set a new direction at Chrysler, which finished its tour through bankruptcy
court Wednesday, completing its deal to join forces with the Italian automaker.

At Ford, a chief executive brought in from the aircraft industry is helping to shake up the company.

But it will be up to the federal government, which will own a majority of General Motors when it
emerges from bankruptcy, to tackle what is perhaps the most difficult challenge in Detroit:
transforming G.M.'s insular culture -- at times as bureaucratic as the government's -- to make the
company more competitive.

If the effort fails, the Treasury may never recoup the \$50 billion it has provided G.M.

``Addressing cultural issues is just as fundamental to our assignment as addressing the balance
sheet or financing,'' said Steven Rattner, the lead adviser to the White House on the automobile
industry.

In just one example, whenever a top G.M. executive was called to appear before lawmakers in
Washington, staff members would prepare a briefing binder as thick as a Manhattan phonebook and hold
multiple meetings to strategize over five minutes of testimony (Fritz Henderson, the new chief
executive, has told employees to stop doing that).

In a Senate hearing Wednesday, Ron Bloom, another adviser on the auto task force, also talked about
the need for G.M. to break longstanding habits that have made the company, with its bloated
structure, lose a step to more nimble competitors.

``General Motors has been kicking problems down the road for a long time,'' Mr.~Bloom said.

Mr.~Rattner and other government officials have repeatedly said they have no interest in running the
company day-to-day. But they are taking a keen interest in shaping the new leadership team.

For example, Mr.~Rattner himself was behind the appointment Tuesday of Edward E.~Whitacre, the
former chief executive at AT\&T, to become G.M.'s new chairman. Four more vacancies on the board
also have to be filled.

The newcomers' primary task will be to make sure that G.M. rewires the corporate DNA that led to a
decades-long slide in market share and tens of billions of dollars in losses.

Mr.~Rattner said he expects no less than the kind of wholesale restructuring that John F.~Welch
Jr.~engineered at General Electric in the 1980s, where he decisively shed businesses, cut costs and
jobs and demanded better performance.

The government, which has said it will not invest any more money in G.M., cannot settle for anything
less, Mr.~Rattner said in an interview.

``We're not going to fail,'' he added. ``This is too important a company.''

If G.M.'s history is any guide, outsiders can have trouble making a difference. For a
quarter-century, a parade of executives and activist investors, including H.~Ross Perot, the
billionaire Kirk Kerkorian and his adviser, the restructuring expert Jerome B.~York, has failed to
shake things up at the highest level.

Those brought into the company to help invigorate it, like Ronald L.~Zarrella, the Bausch \& Lomb
executive who ran its marketing, and the Chicago lawyer Elmer W.~Johnson, who held a top executive
job at G.M., ultimately left the company.

In a famous memo that circulated in Detroit during the late 1980s, Mr.~Johnson complained that the
company was hampered by its inability to execute.

``Teamwork has been replaced by Balkanization,'' he wrote. ``Our culture discourages open, frank
debate among G.M. executives in the pursuit of problem resolution.''

Even within hours of seeking bankruptcy protection on June 1, the company showed old habits that
seemed out of sync with its new, humbled circumstances.

G.M.'s chief financial officer, Ray Young, told reporters that as a ``privately held company,'' G.M.
was not required to regularly disclose details on its financial performance -- even though its new
owners were in fact American and Canadian taxpayers, the United Automobile Workers union and G.M.
bondholders.

After Canadian officials and taxpayer advocates protested, G.M. retreated a day later, promising to
make regular disclosures.

Through the years, outsiders and even company insiders have complained about the months of study
required for simple decisions, like making design changes to a bumper or headlights.

Government officials are vetting new board members who will join Mr.~Whitacre, who admits he has no
car industry experience, in providing oversight of Mr.~Henderson, the successor to Rick Wagoner, who
was forced to resign as chief executive by the Treasury in March.

Mr.~Henderson must now prove that he can make good on his promise to make G.M. leaner and faster.
Government officials decided to give him the title of chief executive, rather than interim chief, to
reassure G.M. troops who were rattled by Mr.~Wagoner's abrupt departure.

Historically, G.M.'s senior leadership group, made up of about 100 managers, met in person four
times a year. Since taking over on March 30, Mr.~Henderson has held three conference calls with the
group, in which he has stressed the need for change and pushed managers to speed up decision-making,
said G.M.'s chief spokesman, Steven J.~Harris.

Mr.~Henderson is leading by example. Upon receiving memos from Mr.~Harris proposing action on
various matters, he told Mr.~Harris not to send similar notes in the future, but instead to provide
an informational copy once a decision has been made.

Mr.~Harris said he disagreed that G.M.'s board had refrained from scrutinizing company activities,
as government officials have implied in public comments.

``They just had a different take on things,'' he said. ``We understand that this is a new
environment and that we're likely to undergo a deeper level of scrutiny.''

Measuring any progress in changing the culture will take time. The results, after all, will be seen
in the new vehicles that the company develops and produces -- and whether they reflect world-class
business practices that are required to win against the best of its global competitors.

``It's not uncomplicated,'' Mr.~Rattner said. ``We hope and believe that it can take place.''


\section{Civic-Minded Chinese Find a Voice Onli}

BEIJING -- There was a time when the story of the 21-year-old waitress who fatally stabbed a
Communist Party official as he tried to force himself on her would have never left the rural byways
of Hubei Province where it took place.

Instead, her arrest last month on suspicion of voluntary manslaughter erupted into an online furor
that turned her into a national hero and reverberated all the way to China's capital, where censors
ordered incendiary comments banned. Local Hubei officials even restricted television coverage and
tried to block travel to the small town where the assault occurred.

On Tuesday, a Hubei court granted the woman, Deng Yujiao, an unexpectedly swift victory, ruling that
she had acted in self-defense and freeing her without criminal penalties.

The case of Ms.~Deng is only the most recent and prominent of several cases in which the Internet
has cracked open a channel for citizens to voice mass displeasure with official conduct,
demonstrating its potential as a catalyst for social change.

The government's reactions have raised questions about how much power officials have to control what
they call ``online mass incidents.'' China's estimated 300 million Internet users, experts say, are
awakening to the idea that, even in authoritarian China, they sometimes can fight City Hall.

``It's about raising the public awareness of democratic ideas -- accountability, transparency,
citizens' rights to participate, that the government should serve the people,'' said Xiao Qiang, a
journalism professor at the University of California, Berkeley, who tracks China's Internet
activity. ``Netizens who are now sharing those more democratic values are using these cases, each
time making inch-by-inch progress.''

China still exerts sweeping and sophisticated control over the Internet, employing thousands of
people to monitor Internet traffic for forbidden material and using software to spot key words that
hint at subversion. But the system is not infallible, and Internet users frequently find ways to
skirt the censors.

Since late last year, online tempests have blown up over a video of an official in Guangdong
Province who assaulted a young girl and bragged that he was above punishment, and a Nanjing city
official who was spotted wearing a \$14,500 Vacheron Constantin watch and smoking \$22-a-pack
cigarettes, evidence of a lifestyle well beyond his means.

Early this year, an online outcry exposed prison officials' cover-up of the beating death of an
inmate. At the moment, outrage is focused on officials in Yunnan Province who battled a rabies
outbreak by dispatching ``killing teams'' that, according to news reports, beat 50,000 dogs to
death.

Not all the crusades are entirely civic-minded. In more than a few cases, virtual mobs have harassed
offending officials, posting personal information and other details. The nickname for such mobs,
``human-flesh search engines,'' hints at their pitiless nature.

But the Internet campaigns have repeatedly produced results. Six officials were punished or fired in
the prison beating. The Nanjing official with the flashy watch was sacked. The Yunnan dog killings
have provoked harsh criticism, even in state-run newspapers.

Most such cases, says Mr.~Xiao, the Berkeley professor, spawn tens or hundreds of thousands of
mentions on Internet blogs and other forums.

But Ms.~Deng's case eclipsed them all, racking up four million posts and counting, he said. Her
story resonates with millions of Chinese who not only are fed up with low-level corruption but also
prize chastity in young women, causes that transcend politics.

``Deng Yujiao is a metaphor for someone who fights back against officials, and of course the
officials are those who spend the taxpayers' money, who are so abusive to ordinary citizens and so
corrupt,'' he said. ``It's almost a stereotype of the online image of officials. That's why this
case becomes so big.''

As she described it to a lawyer, Ms.~Deng was a waitress in a karaoke parlor in rural Badong County,
a Hubei Province backwater along the Yangtze River. Like more than a few such venues, this one
offered ``special services,'' or prostitution, in a backroom spa, the only room with hot water.

On the night of May 10, Ms.~Deng said, she was in the room washing clothes, when a local official,
Huang Weide, came in and demanded that she take a bath with him. She refused, and after a struggle
fled to a bathroom.

But Mr.~Huang and two companions -- including a second official, Deng Guida, who was not related to
Ms.~Deng -- tracked her to the bathroom, then pushed her onto a couch. As they attacked, Ms.~Deng
said, she took a fruit knife from her purse and stabbed wildly. Mr.~Deng fell, mortally wounded.

Ms.~Deng was arrested, investigated for involuntary manslaughter and, after the police reportedly
found pills in her purse, variously described as sleeping pills and antidepressants, sent her to a
mental ward.

But when a blogger, Wu Gan, publicized her case, a cascade of posts crowned her a national hero for
resisting official abuse of power and demanded a fair trial.

Under public pressure, Hubei officials freed her on bail. Mr.~Wu helped recruit a prominent Beijing
law firm to represent Ms.~Deng.

On May 22, Beijing censors ordered Web sites to stop reporting on the case. Four days later,
television and the Internet were cut off in Yesanguan, the town where the attack occurred. The
official explanation for the shutdown was as a ``precaution'' against lightning strikes.

Spurred by the Internet frenzy, Chinese journalists had converged on Badong County. But after
censorship was imposed, local officials began screening outsiders, and some journalists seeking to
report there were beaten. Mr.~Wu's blog was shut down by censors.

Even Yangtze River boat service to Badong was suspended, ostensibly because the docks needed repair,
after protesters vowed to hold a demonstration there.

The two surviving local officials who were involved in the assault have been fired, but no charges
were brought against them.

The ruling on Tuesday, widely reported in state media, was a vindication for Ms.~Deng and her
Internet supporters. But the story may not end there.

Last month, a group of young people abruptly appeared in the middle of downtown Beijing, carrying on
their shoulders a woman wearing a mask and wrapped in white cloth. They laid her on the ground and
arranged signs around her, then took pictures.

The signs read, ``Anyone could be Deng Yujiao.''

The photos immediately appeared on the Internet.

\section{China's Computer Folly}

China has accomplished remarkable things in the past 20 years, including building one of the world's
largest economies. Computers helped speed that development -- and will be even more important in the
future. So Beijing's decision to require that all new personal computers sold in China contain
software that bars access to certain Internet content seems particularly self-destructive and
foolish.

The new rules say all PCs sold in China after June 30 must include special software -- designed by a
company with links to China's military and security agencies -- to filter out pornography and other
``vulgar'' material. Beijing claims that it is trying to protect children. Don't believe it.

In any country, such vague terms would be a frightening license for government intrusion. China's
government, which fears the free flow of ideas, already vigorously restricts Internet content,
including blocking access to Web sites on Tibet, human rights and other politically sensitive
subjects.

Chinese bloggers, dissidents and even some state news media outlets are right to worry that the new
software could be used even more nefariously: to collect personal data and spy on consumer Web
habits.

The contract for the software, meanwhile, was awarded without industry input. There are serious
questions about whether the product will even work.

The last thing China needs is to force the installation of software that could cause millions of
computers to crash. That would feed new resentment against a government already accused of gross
incompetence after thousands of children died in the collapse of shoddily constructed schools in the
2008 earthquake.

International manufacturers probably could force the government to reverse the new rules by
threatening not to sell their products. But they have no history of standing up to Beijing. We hope
they are making a stronger case in private for a rollback than was apparent in the anemic public
statement issued by a coalition of American trade associations. They called for ``an open and
healthy dialogue'' with the government but seemed to go along with the farcical claim that its
intention really was to improve parental control.

If Beijing does not reconsider its foolhardy decision, the new rules would take effect on July 1.
Our advice to Chinese consumers: Buy your PC now.

\section{The Tweet Smell of Success}

Mr.~Rogers had many neighbors, but was always looking to add more. On Facebook, people who have
hundreds of ``friends'' still collect them avidly: ``Will you be my friend, we already have 15 in
common?''

Twitter, the social-networking site of the moment, traffics in a different currency --
``followers,'' who presumably are loyal to you, fascinated by you, enthralled by you. Imagine what
you could accomplish with an army of followers, the lands you could conquer!

Twitter could help you. In the last few months it has plucked a few hundred users from a sea of more
than 30 million and put them on its A-list, deeming them particularly worthy of being followed.

In separating the wheat from the chaff, Twitter has become a kingmaker of sorts, conferring online
stardom to a mix of writers, gadget geeks, political commentators and entrepreneurs.

After being named to the ``suggested user'' list, Twitterers can gain more than 500,000 followers
who get their brief updates via a cellphone or the Internet. A writer with an interest in comic
books can become the expert on comic books; a political pundit with a radio show ends up having a
greater audience online than on the air; and an actor like LaVar Burton, decades away from his glory
days as a star of the TV drama ``Roots,'' has a personal audience of 635,000.

And just as publicity agents used to inveigle syndicated columnists like Walter Winchell into giving
their clients a mention, modern-day publicity hounds are already trying to game the list.

The Web entrepreneur Jason Calacanis declared that he would pay \$250,000 to be on Twitter's
suggestion list for two years. He says the offer was only ``half-real.''

But in an e-mail interview he explained that had Twitter accepted his offer he would have used his
enhanced status to promote his search engine, Mahalo, not himself. ``If they had taken the money I
would have two or three million followers on that account,'' and at 10 cents a follower, he would
consider it to be a bargain.

Twitter is not believed to be profitable, and some wonder if it won't end as a flash in the pan like
other online ventures that could never turn public fervor into a going business. Could selling
placement on the suggested users list be used to generate revenue? ``Not as far as I know,'' said
Jenna Sampson, a spokeswoman for Twitter.

Later in the conversation, though, she hedged a bit. ``Everything is an option at this point,'' she
said, reflecting the company's relatively recent arrival.

Complicating Twitter's decision is the recognition that its users themselves presumably could sell
mentions on their feeds to businesses and personalities. The power of getting placed in the column
of J.~J.~Hunsecker (a thinly disguised version of Winchell), for example, is what drives the plot of
the iconic film ``Sweet Smell of Success.'' It was Hunsecker's wheels that were greased, not the
newspaper's.

Kathleen Hessert, whose company, Sports Media Challenge, advises athletes on how to extend their
popularity online, says that cracking the suggested user list could become an important goal for
some of her clients. Not for clients like Shaquille O'Neal or Peyton Manning, but for the B-list
celebrities, like, say, the solid basketball player Troy Murphy of the Indiana Pacers, who uses
Twitter to share his concern for the environment. (She said that she would like Twitter to put
placement up to a vote that would genuinely represent potential public interest.)

``They are making celebrities by choosing who to follow,'' she said. ``What their system is for
picking people, I am foggy about -- and they want it to be mysterious, I'm thinking.''

Twitter says the list was created in January to solve a vexing problem: people who first subscribe
to the service often are overwhelmed by the experience. Suggested users become a welcome wagon.

``People were signing up and then they weren't following anyone,'' said Biz Stone, one of Twitter's
co-founders. Did he realize he was helping to create an arbiter of popularity? ``We didn't think
that far ahead,'' he said.

The list is cobbled together by a team of employees whose identities were withheld, lest they be
bombarded with Twitterers trying to manipulate the process. The company says it compiles the list by
tracking fast-growing accounts and then deciding whom to anoint.

Ms.~Sampson said ``there's sort of a criteria'' for the list ``but not really.'' Twitter says it
wants to highlight personally revealing accounts, not promotional ones, although businesses like The
New York Times, JetBlue and the N.F.L. are represented. In essence, the list indicates ``this is how
we think Twitter should be used,'' she said.

Web users are already well aware of the role that powerful sites like Apple's iTunes, Netflix, and
Amazon.com play in the promotion of music, films and books. Twitter is taking this influence in a
new direction, however, by applying its recommendation power to people instead of products.

Twitter's list includes many celebrities as well, like Lance Armstrong, Mr.~O'Neal and Britney
Spears. But, naturally, it is the once-obscure users, like the science technology writer Steven
Johnson, who are most changed by the experience.

Twitter doesn't notify the people on its list; Mr.~Johnson says he only learned about it in January
when he was on a book tour in the Northwest and suddenly received 200 e-mail messages in a matter of
hours, each informing him that he had a new follower. ``I sent out a tweet saying 'where did you all
come from?' '' he recalled.

In short order, Mr.~Johnson's following rose from 5,000 to more than 630,000. (And the number keeps
growing every hour.) ``That is a big force,'' he said. ``I'm not Oprah, but I can get some chatter
about something going pretty effectively -- I put a link to something and I can see it repeated.''

In the process, Mr.~Johnson said, he has been able to witness the rising importance of new-media
outlets and the lessening influence of traditional media outlets like Time magazine, which recently
printed his essay on (surprise!) the transformative power of Twitter.

That essay was featured on the cover with a sample tweet of Mr.~Johnson's. ``It's funny, everybody
has been asking me, you got your Twitter ID on the cover of Time magazine, you must be getting an
insane amount of followers,'' he said. ``And I say it's nothing compared to the steady influx you
get from being on the suggested user list.''

\section{Real-Time Criticism of CNN's Iran Coverage}

Cable news normally serves as the front line for breaking news, but the channels largely took the
weekend off as Tehran exploded in protests after Iran's presidential election.

The performance of the American cable news, especially CNN, spawned an online protest by thousands
on Saturday and Sunday, showing that viewers can try to pressure news organizations about their
coverage in real time via the Internet. Fox News Channel and MSNBC also were said to have covered
the protests in limited ways.

Protesters' comments on Twitter were quickly noticed by CNN, which defended its coverage. The social
networking blog Mashable said Twitter was acting as a ``media watchdog.''

CNN had reports from Tehran throughout Saturday, including some from Christiane Amanpour, its chief
international correspondent. But it did not provide the kind of wall-to-wall coverage that some had
expected.

It was a departure for CNN, known for its breaking news coverage, including its celebrated reporting
during the Tiananmen Square crackdown 20 years ago. But the Tehran protests were not covered with
rolling live coverage for hours at a time.

Untold thousands used the label ``CNNfail'' on Twitter to vent their frustrations. Steve LaBate, an
Atlanta resident, said on Twitter, ``Why aren't you covering this with everything you've got?''
About the same time, CNN was showing a repeat of Larry King's interview of the stars of the
``American Chopper'' show. For a time, new criticisms were being added on Twitter at least once a
second.

Andrew Sullivan, a blogger for The Atlantic, wrote, ``There's a reason the MSM is in trouble,''
using the blogosphere abbreviation for mainstream media.

CNN said, ``We share people's expectations of CNN and have delivered far more coverage of the
Iranian election and aftermath than any other network.''

Journalists in Tehran were working in difficult and sometimes dangerous conditions. The BBC said its
correspondent John Simpson and a camera operator had been briefly arrested after filming in the
streets. Jim Sciutto, an ABC News correspondent in Tehran, said that the police had confiscated a
camera and footage. ``We are shooting protests and police violence on our cellphones,'' he wrote on
Twitter.

Some Americans relied on British networks. A report from Channel 4 was spread widely on the
Internet. In the video, the correspondent Lindsey Hilsum said of Iran, ``I feel like I went to sleep
in one country and woke up in another.'' BRIAN STELTER

\section{U.S.~Steps Gingerly Into Tumult in Iran}

WASHINGTON -- The Obama administration says it has tried to avoid words or deeds that could be
portrayed as American meddling in Iran's presidential election and its tumultuous aftermath.

Yet on Monday afternoon, a 27-year-old State Department official, Jared Cohen, e-mailed the
social-networking site Twitter with an unusual request: delay scheduled maintenance of its global
network, which would have cut off service while Iranians were using Twitter to swap information and
inform the outside world about the mushrooming protests around Tehran.

The request, made to a Twitter co-founder, Jack Dorsey, is yet another new-media milestone: the
recognition by the United States government that an Internet blogging service that did not exist
four years ago has the potential to change history in an ancient Islamic country.

``This was just a call to say: 'It appears Twitter is playing an important role at a crucial time in
Iran. Could you keep it going?' '' said P.J.~Crowley, the assistant secretary of state for public
affairs.

Twitter complied with the request, saying in a blog post on Monday that it put off the upgrade until
late Tuesday afternoon -- 1:30 a.m. Wednesday in Tehran -- because its partners recognized ``the
role Twitter is currently playing as an important communication tool in Iran.'' The network was
working normally again by Tuesday evening.

The State Department said its request did not amount to meddling. Mr.~Cohen, they noted, did not
contact Twitter until three days after the vote was held and well after the protests had begun.

``This is completely consistent with our national policy,'' Mr.~Crowley said. ``We are proponents of
freedom of expression. Information should be used as a way to promote freedom of expression.''

The episode demonstrates the extent to which the administration views social networking as a new
arrow in its diplomatic quiver. Secretary of State Hillary Rodham Clinton talks regularly about the
power of e-diplomacy, particularly in places where the mass media are repressed.

Mr.~Cohen, a Stanford University graduate who is the youngest member of the State Department's
policy planning staff, has been working with Twitter, YouTube, Facebook and other services to
harness their reach for diplomatic initiatives in Iraq and elsewhere.

Last month, he organized a visit to Baghdad by Mr.~Dorsey and other executives from Silicon Valley
and New York's equivalent, Silicon Alley. They met with Iraq's deputy prime minister to discuss how
to rebuild the country's information network and to sell the virtues of Twitter.

Referring to Mir Hussein Moussavi, the main Iranian opposition candidate, Mr.~Crowley said, ``We
watched closely how Moussavi has used Facebook to keep his supporters informed of his activities.''

Tehran has been buzzing with tweets, the posts of Twitter subscribers, sharing news on rallies,
police crackdowns on protesters, and analysis of how the White House is responding to the drama.

With the authorities blocking text-messaging on cellphones, Twitter has become a handy alternative
for information-hungry Iranians. While Iran has also tried to block Twitter posts, Iranians are
skilled at using proxy sites or other methods to circumvent the official barriers.

A Twitter account called IranNewsNow sent a message to CNN's Twitter account that read, ``don't
listen to what iran gov says u can or can't do! You can report the pics/vids coming from Twitter!''

An account called StopAhmadi wrote on Tuesday evening, ``We need ppl around world helping to raise
the issues put pressure on Iranian gvmt.'' It posted links to pictures from Tehran, including one
that showed a man bleeding profusely from his chest, surrounded by protesters.

There were also suspicions that some pro-government forces might be using new-media outlets to send
out misinformation. One popular opposition site, Persiankiwi, warned its followers on Tuesday to
ignore instructions from people with no record of reliable posts.

In addition to Twitter, YouTube has been a critical tool to spread videos from Iran when traditional
media outlets have had difficulty filming the protests or the ensuing crackdown. One YouTube
account, bearing the user name ``wwwiranbefreecom,'' showed disturbing images of police officers
beating people in the streets. On Monday, Lara Setrakian, an ABC News journalist, put out a call for
video on Twitter, writing, ``Please send footage we can't reach!''

The BBC's Persian-language television channel said that for a time on Tuesday, it was receiving
about five videos a minute from amateurs, even though the channel is largely blocked within Iran.
One showed pro-government militia members firing weapons at a rally.

``We've been struck by the amount of video and eyewitness testimony,'' said Jon Williams, the BBC
world news editor. ``The days when regimes can control the flow of information are over.''

As new media proliferate, however, traditional journalists are having a harder and harder time.

Journalists were told on Tuesday that they could not cover protests without permission. The
restrictions ``effectively confine journalists to their offices,'' a spokesman for the BBC said.

Still, many ventured out into the streets to witness pro- and antigovernment protests, at
considerable risk. At the Laleh Hotel in central Tehran, the Time magazine columnist Joe Klein said,
``A number of journalists were coming back from the streets pretty badly beaten.''

As their visas expired, journalists were looking for any chance to report. Jim Marshall, the last
Sky News staff member in Tehran, was barred from reporting, so he went shopping instead and came
upon thousands of supporters of President Mahmoud Ahmadinejad at a rally.

``I kept shopping, and they kept demonstrating,'' he wrote in a blog post. ``This was turning into a
test of wills. How much longer could I shop without slipping into reporting?'' Once he realized he
was carrying a notepad in his pocket, he swiftly left the scene.

\section{China's College Entry Test Is an Obsession}

TIANJIN, China -- For the past year, Liu Qichao has focused on one thing, and only one thing: the
gao kao, or the high test.

Fourteen to 16 hours a day, he studied for the college entrance examination, which this year will
determine the fate of more than 10 million Chinese students. He took one day off every three weeks.

He was still carrying his textbook from room to room last Sunday morning before leaving for the exam
site, still reviewing materials during the lunch break, still hard at work Sunday night, preparing
for Part 2 of the exam that Monday.

``I want to study until the last minute,'' he said. ``I really hope to be successful.''

China may be changing at head-twirling speed, but the ritual of the gao kao (pronounced gow kow)
remains as immutable as chopsticks. One Chinese saying compares the exam to a stampede of ``a
thousand soldiers and 10 horses across a single log bridge.''

The Chinese test is in some ways like the American SAT, except that it lasts more than twice as
long. The nine-hour test is offered just once a year and is the sole determinant for admission to
virtually all Chinese colleges and universities. About three in five students make the cut.

Families pull out all the stops to optimize their children's scores. In Sichuan Province in
southwestern China, students studied in a hospital, hooked up to oxygen containers, in hopes of
improving their concentration.

Some girls take contraceptives so they will not get their periods during the exam. Some well-off
parents dangle the promise of fabulous rewards for offspring whose scores get them into a top-ranked
university: parties, 100,000 renminbi in cash, or about \$14,600, or better.

``My father even promised me, if I get into a college like Nankai University in Tianjin, 'I'll give
you a prize, an Audi,' '' said Chen Qiong, a 17-year-old girl taking the exam in Beijing.

Outside the exam sites, parents keep vigil for hours, as anxious as husbands waiting for their wives
to give birth. A tardy arrival is disastrous. One student who arrived four minutes late in 2007 was
turned away, even though she and her mother knelt before the exam proctor, begging for leniency.

Cheating is increasingly sophisticated. One group of parents last year outfitted their children with
tiny earpieces, persuaded a teacher to fax them the questions and then transmitted the answers by
cellphone. Another father equipped a student with a miniscanner and had nine teachers on standby to
provide the answers. In all, 2,645 cheaters were caught last year.

Critics complain that the gao kao illustrates the flaws in an education system that stresses
memorization over independent thinking and creativity. Educators also say that rural students are at
a disadvantage and that the quality of higher education has been sacrificed for quantity.

But the national obsession with the test also indicates progress. Despite a slight drop in
registration this year -- the first decline in seven years -- five million more students signed up
for the test than did so in 2002.

China now has more than 1,900 institutions of higher learning, nearly double the number in 2000.
Close to 19 million students are enrolled, a sixfold jump in one decade.

Liu Qichao, 19, a big-boned student with careful habits, plans to be the first in his family to go
to college. ``There just were not a lot of universities then,'' said his father, Liu Jie, who
graduated from high school in 1980 and sells textile machinery. His son harbors hopes of getting
into one of China's top universities.

But the whole family was shaken by the results of his first try at the gao kao last June.

The night before the exam, he lingered at his parents' bedside, unable to sleep for hours. ``I was
so nervous during the exam my mind went blank,'' he said. He scored 432 points out of a possible
750, too low to be admitted even to a second-tier institution.

Silence reigned in the house for days afterward. ``My mother was very angry,'' he said. ``She said,
'All these years of raising you and washing your clothes and cooking for you, and you earn such a
bad score.'

``I cried for half a month.''

Then the family arrived at a new plan: He would enroll in a military-style boarding school in
Tianjin, devoting himself exclusively to test preparation, and retake the test this June.

Despite the annual school fee of 38,500 renminbi (about \$5,640) -- well above the average annual
income for a Chinese family -- he had plenty of company.

One of his classmates, Li Yiran, a cheerful 18-year-old, estimated that more than one-fourth of the
seniors at their secondary school, Yangcun No.~1 Middle School, were ``restudy'' students.

Ms.~Li said she learned the hard way about the school's strict regimen. When her cellphone rang in
class one day, the teacher smashed it against the radiator. Classes continue for three weeks
straight, barely interrupted by a one-day break.

Days after most of their classmates left for home, Mr.~Liu and Ms.~Li were still holed up last week
in their classrooms. Mr.~Liu's wrist was bruised from pressing the edge of his blue metal desk,
piled with a foot-high stack of textbooks.

Ms.~Li's breakfast was a favorite among test-takers: a bread stick next to two eggs, symbolizing a
100 percent score.

Hours after they finished the test on Monday, both students had collected the answers from the
district education bureau and begun the laborious process, with the help of their teachers, of
estimating their scores.

Mr.~Liu calculated that his score leaped by more than 100 points over last year's dismal
performance. But he was still downcast, uncertain whether he would make the cutoff to apply to
top-tier universities. The cutoff mark can vary by an applicant's place of residence and ethnicity.

Ms.~Li, on the other hand, was exhilarated by her estimate of 482.5, figuring it was probably high
enough for admittance to a college of the second rank.

By Wednesday evening, both were buoyed by news of the cutoff scores for their district. His
estimated mark was well above the one needed to apply to first-tier schools, and hers was a solid
five points above the notch for the second tier.

Before the test, Ms.~Li's aunt warned her that this was her last chance for a college degree. Even
if she knelt before her mother and begged, her aunt said, her mother would refuse to let her take
the test again.

But Ms.~Li, a hardened veteran of not one but two gao kao ordeals, had a ready retort: ``Come on.
Even if my mother kneels down before me, I will refuse to take this test again.''

\section{China Orders Patches to Planned Web Filter}

BEIJING -- A designer of censorship software that the Chinese government requires to be preinstalled
on computers sold in China has been ordered to fix potential security breaches in the software, the
newspaper China Daily reported Monday. The report was an indication that the government still
supports use of the software despite heated debate over it.

The software, called Green Dam-Youth Escort, has come under attack by many computer users in China
for both political and technical reasons. Critics say that although the Chinese government insists
that the software will be used only to block access to pornography Web sites, the software's actual
use will be to block any site with content deemed politically objectionable, like the Tibet issue or
the 1989 Tiananmen killings.

The government says all computers sold in China must have the software installed by July 1.

Some computer experts who have studied the software said last week that it was so flawed that it
could allow hackers to monitor a user's Internet activity, steal personal data or plant viruses. One
expert, J.~Alex Halderman, a computer science professor at the University of Michigan, has posted on
the Internet a report on Green Dam's vulnerabilities.

Rather than agreeing to scrap the software altogether, the Chinese government has responded to the
technical criticisms by ordering that the potential security breaches be eliminated.

``The Ministry of Industry and Information Technology told us to make the software safer as soon as
a series of security vulnerabilities were found,'' Zhang Chenmin, the general manager of Jinhui
Computer System Engineering, which helped design the software, told China Daily.

Mr.~Zhang admitted that the software had systemic flaws that would allow hackers to attack computers
that used the program, ``just like any other software of this type.''

Mr.~Halderman said in an interview last week that it had only taken a few hours for him and his
students to infiltrate a computer loaded with Green Dam and force it to crash. A skilled hacker
could take over the computer to mine personal data or hitch it to other infected machines in a
malevolent network known as a botnet, he added.

Debate about the software exploded on the Internet last week as Chinese Internet users increasingly
learned about the software requirement issued by the ministry. China Daily reported Monday that
surveys done by four of China's most popular Web portals showed that four in five Internet users
would not use the software or would have it uninstalled.

Early reports had indicated that the government might simply require Green Dam to be included on a
CD packaged with new computers, so users would have the option to install it. But it became apparent
last week that the government was insisting that all computer makers preinstall the software by July
1. Foreign computer makers learned of the directive just three weeks ago and have been asking the
Chinese government to reconsider the rules.

Meanwhile, Solid Oak Software, a company based in Santa Barbara, Calif., has accused the designers
of Green Dam of stealing programming code from software called Cybersitter that was developed by
Solid Oak. Cybersitter is designed to block Web sites deemed to be pornographic, violent or
offensive. Solid Oak says that some of the ``blacklist'' files used by Green Dam were originally
developed for Cybersitter.

``I cannot deny that the two filters' databases of blacklisted URL addresses might share
similarities,'' Mr.~Zhang said in China Daily. ``After all, they are all well-known international
pornographic Web sites that all porn filters are meant to block. But we didn't steal their
programming code.''






\section{China Holds Firm on Software Filter, U.S. Firms Say}

BEIJING -- American computer makers say the Chinese government has not backed down from a
requirement that Internet censorship software be preinstalled on all computers sold in China after
July 1, despite reports this week that the rule had been relaxed.

In a further sign that Chinese officials are trying to assert more Internet control, the city of
Beijing wants to recruit 10,000 volunteers by summer's end to monitor online content, said an
employee of the city government's Spiritual Civilization Office.

The plan was presented in a document submitted Tuesday by the Beijing Internet Administration Office
during a meeting in which city officials discussed ``purifying social civilization,'' said the
employee, who identified herself only by her surname, Guo. She said she had no further details on
the plan.

Chinese authorities have also sought to assert control by directly warning some online services. On
Thursday, for example, a government-supported Internet watchdog group criticized Google's
Chinese-language Web site for linking to ``pornographic and vulgar'' sites. The group, the China
Internet Illegal Information Reporting Center, said www.google.cn had already been criticized in
January and April and that it must purge the offending links.

China's central and local governments use a vast array of programs and human monitors to block
Internet content deemed pornographic or politically harmful to the Communist Party, like Web sites
discussing Tibet or the Falun Gong, the banned spiritual movement. The system of censorship is
nicknamed the Great Firewall, and savvy computer users in China usually use software to circumvent
it, thinking little of it. But the government's new rule that computer makers install censorship
software on computers has inflamed antigovernment sentiment among Chinese computer users.

Many people say the software, called Green Dam-Youth Escort, will be used to block Web sites with
politically unacceptable content, even though officials insist that the software will be used
primarily to censor pornography.

Computer experts also discovered severe weaknesses in the software that would let hackers hijack
computers. Chinese officials say they have ordered the developers to fix these problems. Trade
groups representing major American computer makers, including Hewlett-Packard and Dell, which have
significant market shares in China, have been asking the Chinese government to rescind the
requirement that Green Dam be preinstalled on computers but have seen no change in the Chinese
position.

Four trade groups based in the United States have sent a statement to the Chinese government asking
it to ``reconsider implementing its new mandatory filtering software requirement.''

On Wednesday, the major American computer makers said they had yet to hear anything concrete from
China regarding the possibility of making installations of Green Dam optional.

Lenovo, the largest computer maker in China, did not respond to requests for comment.

Confusion was sown Monday when China Daily, the country's official English-language newspaper,
quoted an unnamed official in the software department of the Ministry of Industry and Information
Technology saying that the government was requiring that the software be offered on a CD-ROM
packaged with new computers or be placed on hard drives as setup files only.

Some people interpreted that as the government backing down from its rule that Green Dam be
installed, but soon it became clear that the official was not speaking in an authoritative role.
Employees at the ministry's software department, reached by telephone, refused to clarify the
government's position. No government official has given any statement this week indicating that the policy has changed.

On May 19, the ministry issued a directive to computer makers that required the preinstallation of Green Dam on hard drives or a CD-ROM with installation software that usually is packaged with computers. The directive also requires that Green Dam be saved in backup files.

The directive makes clear that the government intends to ensure universal use of Green Dam on new computers in China.

The software and its two Chinese developers are coming under legal attack by Solid Oak Software, a
company based in Santa Barbara, Calif., which is accusing the developers, Jinhui Computer System
Engineering and Dazheng Human Language Technology, of stealing programming codes from software that
Solid Oak developed.

Solid Oak said it had asked Hewlett-Packard and Dell to halt distribution of any computers with
Green Dam software.

\section{Apple Fills in Some Gaps With Latest iPhone}

Assessing the 2007 and 2008 iPhone models was an excruciating experience. You were torn in half --
between your heart and your head.

Your emotions were swept away by everything Apple does so well: beauty, polish, elegance, simplicity
and the thrill of interaction. (Those were not, ahem, phrases typically used to describe existing
cellphones.)

Meanwhile, your brain kept waving its little hand in the back of the classroom. ``But the camera's
terrible!'' it would say. ``It can't record video! There's no voice dialing! No copy and paste! The
iPhone can't even send picture messages -- even \$20 starter phones can do that!''

But 21 million iPhone sales later, it's become clear that the heart usually manages to shut the head
up.

With the iPhone 3GS, in stores Friday, Apple is finally throwing your head a crumb. After two
years, the iPhone's designers have finally gotten over whatever weird objections they had to
providing those basic functions.

Better yet, Apple intends to give many of those features, and dozens more, to everyone who has ever
bought an iPhone.

If you do buy the iPhone 3GS, you get twice the storage -- 16 gigabytes -- for the same \$200 price
as before. For \$300, you can even buy a 32-gigabyte model, enough to hold the entire ``Lord of the
Rings'' trilogy, the DVD extras and 75 gazillion songs.

(These are prices for new customers. If you bought last year's iPhone relatively recently, you'll
have to pay \$200 extra for the new one, a point of outrage among the Apple faithful. Unfortunately,
that's just the way the subsidized-cellphone business works. On the other hand AT\&T has just
announced a special offer: if your existing iPhone will be ``upgrade eligible'' this July, August or
September, you can get the new iPhone now, for the new-customer price. More on this complex subject
is at nytimes.com/pogue.)

You can still buy last year's model, the iPhone 3G, for \$100. But do find a way to afford the new
one. It looks identical to last year's iPhone, but its faster circuitry makes a huge difference.
(The S stands for speed, says Apple.) If you're used to the old iPhone, the speed boost hits you
between the eyes, especially when you're opening programs, playing games and loading Web pages.

The built-in three-megapixel camera is much better, too. The camera still tends to blur moving
subjects, and even still lifes aren't as crisp as from an actual camera. But the color and clarity
are definitely improved, especially in low light.

The new autofocus feature lets you tap the screen preview at the spot where you want the exposure,
white balance and focus to be calculated. Except when the subject is a few inches away, you don't
see much difference in the focusing -- but your tap location can make a big difference in the
brightness and color (exposure and white balance) of the finished photo. (You can see sample photos
at nytimes.com/personaltech.)

Better yet, the 3GS now captures video. It's the real deal: sharp, smooth, 30 frames a second. Once
again, it's not quite what you'd get from a proper digital camera or a Flip camcorder--it tends to
``blow out'' the bright areas -- but it's darned close.

You can't beat the capacity, either; in theory, the 32-gigabyte iPhone can capture 17 hours of video
-- just enough for the elementary-school talent show.

With a fingertip, you can trim the ends of a captured video and then upload it to YouTube or
MobileMe, right from the phone. (That part, it does much better than a digital camera.)

The new voice-control feature may be the most useful change of all. Hold down the iPhone's Home
button for a moment, say ``Call mom's cell'' or ``Dial 800-555-1212,'' and the iPhone places your
call, crisply and accurately. (Yeah, I know: welcome to 2003.) This feature goes a long way toward
addressing what's always been the iPhone's weakest feature: the number of steps required to place a
call.

The iPhone also recognizes spoken iPod commands like ``Play songs by Abba'' or ``What song is
this?''

The new Compass program looks like a classier version of a regular Cub Scout compass -- great when
you emerge, disoriented, from the subway. In Google Maps, it adds an indicator beam, showing which
way you are facing on the map. No longer must you walk in a circle, staring at the iPhone map like
an idiot, just to figure out which way is up.

The iPhone 3GS also gains what Apple calls an oleophobic screen. It may sound like an irrational
fear of yodelers, but in fact, it's a coating that lets you wipe away fingerprints with a single rub
on your clothes. It really works to keep the iPhone looking new longer. Maybe fewer people will now
bury the iPhone's gorgeous, slim shape in a homely, bulky case.

Finally, the iPhone 3GS harbors a better, beefier battery, thereby confronting another chronic
complaint. It gives you about 25 percent more life a charge (five hours talk time or 30 hours of
music), easily enough to last at least a day of moderate use. As Palm Pre owners know, that's rare
on a 3G superphone.

There are dozens more new features on the iPhone 3GS -- but the really exciting part is that older
iPhones can get them, too. They're part of a free software upgrade called iPhone 3.0. (You get the
upgrade when you sync your phone to iTunes. For \$10, the iPod Touch can get this upgrade, too.)

Chief among them: the long-awaited copy and paste commands, which appear at your fingertips when you
double-tap text in most programs. Now you can paste text and graphics from a Web site into an e-mail
message, for example, or copy an address from a text message into your calendar.

There's Bluetooth stereo audio, too, meaning that you can listen to your music with cordless
headphones, leaving the iPhone itself in your pocket or backpack. (It's available on the iPhone 3G
S, 3G and the current iPod Touch.) A handy voice-recording app comes complete with trim editing and
e-mailing commands, thereby turning your iPhone into a high-quality, huge-capacity digital audio
recorder.

If you have a MobileMe account (\$100 a year), you can also make your iPhone beep for two minutes --
and display a plaintive message on the screen -- when you've misplaced it. How many times have you
wished your cellphone had that feature?

The 3.0 software also brings, at last, picture and video messaging (known as MMS) to the iPhone 3G
and 3GS -- or it will, once AT\&T turns on this feature later this summer.

The iPhone app store offers a staggering 50,000 instantly downloadable programs, in every
conceivable category; it's become a crucial reason, maybe the crucial reason, to get an iPhone in
the first place. These programs are getting very sophisticated indeed.

Documents to Go (\$5) lets you create and edit Word documents right on the iPhone, for example;
programs like Gokivo and TomTom will bring real, spoken, turn-by-turn GPS navigation to the iPhone.

But the more fun you have trying out these apps, the more desperately you need a way to manage them.
After all, the iPhone can now hold 176 apps on 11 side-by-side Home screens.

Therefore, the new universal-search feature could not arrive at a better time. Type a few
characters, and up pops a list of every match in your calendar, address book, notes, music stash,
e-mail headers (subject, to/from, and first few body lines) -- and apps. That's right: you can now
call up a program by typing a bit of its name.

All of these changes make it much harder to resist the iPhone on intellectual, feature-counting
grounds. The new iPhone doesn't just catch up to its rivals -- it vaults a year ahead of them.

At this point, the usual 10 rational objections to the iPhone have been whittled down to about
three: no physical keyboard, no way to swap the battery yourself and no way to avoid using AT\&T as
your carrier.

In short, the substantially improved, still elegant iPhone 3GS makes it dangerously easy for your
heart and your head to agree.

\section{Iranian Rally Is Dispersed as Voting Errors Are Admitted}

TEHRAN -- Hours after a warning from the powerful Revolutionary Guards not to return to the streets,
about a thousand protesters defiantly gathered in central Tehran on Monday and were quickly
dispersed in an overwhelming show of force by police who used clubs and tear gas.

The protesters, far fewer than the numbers who had attended mass rallies last week, turned out
despite the warning, on the Guards' Web site, that they would face a ``revolutionary confrontation''
if they continued to challenge the results of the June 12 election and their country's supreme
leader, who has pronounced the ballot to be fair.

Security forces descended on a separate crowd of hundreds who had gathered at a square to
memorialize Neda Agha-Soltan, the young woman whose shooting death at a rally over the weekend was
captured on video and circulated around the world.

Despite the Ayatollah's insistence of the election's legitimacy, Iran's most senior panel of
election monitors, in the most sweeping acknowledgment that the election was flawed, said Monday
that the number of votes cast in 50 cities exceeded the actual number of voters, according to a
state television report.

The discrepancies could affect some three million ballots of what the government says was 40 million
cast, giving the official victory to the incumbent president, Mahmoud Ahmadinejad.

The authorities insisted that the discrepancies did not violate Iranian law. The Guardian Council,
charged with certifying the election, said it was not clear whether they would decisively change the
result, which placed Mir Hussein Moussavi -- who contends the election was stolen from him -- in a
distant second.

He has urged his supporters to continue their defiance, but he could face arrest for doing so.

``Moussavi's calling for illegal protests and issuing provocative statements have been a source of
recent unrests in Iran,'' Ali Shahrokhi, head of parliament's judiciary committee, semi-official
Fars news agency reported, according to Reuters. ``Such criminal acts should be confronted firmly.''

He added: ``The ground is paved to legally chase Moussavi.''

Mr.~Moussavi, the more moderate of the candidates, used a posting on his Web site Sunday night to
urge his supporters to demonstrate peacefully, despite warnings from Iran's supreme leader,
Ayatollah Ali Khamenei, that no protests of the vote would be allowed.

``Protesting to lies and fraud is your right,'' Mr.~Moussavi said.

In an apparent response, a Guards statement Monday told protesters to ``be prepared for a resolution
and revolutionary confrontation with the Guards, Basij and other security forces and disciplinary
forces'' if they took their protests into a second week, news reports said.

The Basij is a militia accused by the protesters of brutally repressing demonstrations that
culminated in a day of bloodshed on Saturday that ended in the deaths of at least 10 protesters,
according to the state television.

The Guards told demonstrators Monday to ``end the sabotage and rioting activities,'' calling their
protests a ``conspiracy'' against Iran. The warning echoed remarks by a Foreign Ministry spokesman
who blamed western governments and media for the unrest.

The official result gave Mr.~Ahmadinejad 63 percent of the ballots -- an 11-million vote advantage
-- to Mr.~Moussavi's 34 percent. Turnout was put at 85 percent.

At a news conference Monday, Hassan Qashqavi, a Foreign Ministry spokesman, called the turnout a
``brilliant gem which is shining on the peak of dignity of the Iranian nation.''

He accused unidentified western powers and news organizations, which are operating under extremely
tight official restrictions, of spreading unacceptable ``anarchy and vandalism.'' But, he said, the
outcome of the vote would not be changed. ``We will not allow western media to turn this gem into a
worthless stone,'' he said.

Mr.~Qashqavi drew comparisons with American election results.

``No one encouraged the American people to stage a riot'' because they disagreed with the
re-election of George W.~Bush in 2004, he said. Britain's Foreign Office said Monday that because of
the continuing unrest it would evacuate the families of staff members based in Iran. A spokeswoman,
who spoke in return for anonymity under civil service rules, said the violence in Tehran ``had a
significant impact on the families of our staff who have been unable to carry on their lives as
normal.''

Quoted earlier by Press TV, Abbas Ali Kadkhodaei, the spokesman for the 12-member Guardian Council
denied claims by another losing candidate, Mohsen Rezai, that irregularities had occurred in up to
170 voting districts.

``Statistics provided by the candidates, who claim more than 100 percent of those eligible have cast
their ballot in 80 to 170 cities are not accurate -- the incident has happened in only 50 cities,''
Mr.~Kadkhodaei said.

But he said that a voter turnout in excess of the registered voting list was a ``normal phenomenon''
because people could legally vote in areas other than those in which they were registered.
Nonetheless, some analysts in Tehran said, the number of people said to be traveling on election day
seemed unusually high.

The news emerged on the English-language Press TV Web site late Sunday as a bitter rift among Iran's
ruling clerics deepened. As increasingly violent protests have swirled through Tehran since the
elections, Ayatollah Khamenei has ordered the Guardian Council to investigate the opposition's
allegations of electoral fraud. The council itself has offered a random partial recount of 10
percent of the ballot.

Mr.~Kadkhodaei said the Guardian Council could recount votes in areas where irregularities were said
by the opposition to have occurred. But ``it has yet to be determined whether the possible change in
the tally is decisive in the election results.''

The opposition has alleged a total of 646 electoral irregularities and is demanding that the vote be
annulled. But in a sermon at Friday prayers last week Ayatollah Khamenei mocked the idea that the
huge margin attributed to Mr.~Ahmadinejad could have been won through fraud.

On Sunday, the police detained five relatives of Ali Akbar Hashemi Rafsanjani, a former president
who leads two influential councils and openly supported Mr.~Moussavi's election. The relatives,
including Mr.~Rafsanjani's daughter, Faezeh Hashemi, were released after several hours.

The developments, coming one day after protests here in the capital and elsewhere were crushed by
police officers and militia members using guns, clubs, tear gas and water cannons, suggested that
Ayatollah Khamenei was facing entrenched resistance among some members of the elite.

Though rivalries have been part of Iranian politics since the 1979 revolution, analysts said that
open factional competition amid a major political crisis could hinder Ayatollah Khamenei's ability
to restore order.

There was no verifiable accounting of the death toll from the bloodshed on Saturday, partly because
the government has imposed severe restrictions on news coverage and warned foreign reporters who
remained in the country to stay off the streets.

It also ordered the BBC's longtime correspondent expelled and Newsweek's correspondent detained.

State television said that 10 people had died in the weekend clashes, while radio reports said 19.
The news agency ISNA said 457 people had been arrested.

In the network of Internet postings and Twitter messages that has become the opposition's major tool
for organizing and sharing information, a powerful and vivid new image emerged: a video posted on
several Web sites that showed a young woman, called Neda, her face covered in blood. Text posted
with the video said she had been shot. It was not possible to verify the authenticity of the video.

The Web site of another reformist candidate, Mehdi Karoubi, referred to her as a martyr who did not
``have a weapon in her soft hands or a grenade in her pocket but became a victim by thugs who are
supported by a horrifying security apparatus.''

On Monday,several hundred gathered at Haft-e-Tir Square to attend a memorial for the woman, but were
quickly chased away by police and security officers. One witness said the officers violently beat
and arrested people who attended the vigil, then attacked people nearby who were seen filming the
events from their homes.

Mr.~Moussavi, meanwhile, was not seen in public on Sunday but showed no sign of yielding. In his Web
posting, he urged followers to ``avoid violence in your protest and behave as though you are the
parents that have to tolerate your children's misbehavior at the security forces.''

He also warned the government to ``avoid mass arrests, which will only create distance between
society and the security forces.''

\section{Despite Law, Job Conditions Worsen in China}

DONGGUAN, China -- Liu Pan, a 17-year-old factory worker, was crushed to death last April when the
machine he was operating malfunctioned.

Somehow Mr.~Liu became stuck in the machine, his sister Liu Yan recalled during a tearful interview
in a village near the factory.

``When we got his body, his whole head was crushed,'' Ms.~Liu said. ``We couldn't even see his
eyes.''

Investigating the accident, inspectors found a series of labor and safety violations at the factory,
Yiuwah Stationery, which supplies cards, gift boxes and other paper goods to Disney, the British
supermarket chain Tesco and other companies.

The investigators also discovered that Mr.~Liu was hired illegally, at 15, below the legal age limit
here. Disney has called the situation at the factory ``unacceptable.''

In a statement issued Wednesday, Disney said it had instructed its vendors and licensees to ``cease
new orders of any Disney-branded products in the Yiuwah factory'' until conditions were improved.

A spokesman for Tesco said that company was also working to improve conditions at the factory.

While the accident at the Yiuwah factory was particularly tragic, working conditions elsewhere are
worsening. A year and a half after a landmark labor law took effect in China, experts say conditions
have actually deteriorated in southern China's export-oriented factories, which produce many of
America's less expensive retail goods.

With China's exports reeling and unemployment rising because of the global slowdown, there is
growing evidence that factories are ignoring or evading the new law, and that the government is
reluctant to enforce it.

Government critics say authorities fear that a crackdown on violators could lead to mass layoffs and
even social unrest.

``The economic downturn has given regulators the perfect excuse to ignore the law,'' says Zhang
Zhiru, director of the Shenzhen Chunfeng Labor Dispute Service, a nonprofit group that supports
workers. ``I don't see any fundamental change.''

But workers are fighting back. Earlier this month, the government said Chinese courts were trying to
cope with a soaring number of labor disputes, apparently from workers emboldened by the promise of
the new contract labor law.

The number of labor disputes in China doubled to 693,000 in 2008, the first year the law was in
effect, and are rising sharply this year, the government says.

The law requires that all employees have a written contract that complies with minimum wage and
safety requirements. It also strengthens the monopoly state-run labor union and makes it more
difficult for companies to use temporary workers or to dismiss employees.

Western companies that outsource to China say they have stepped up their monitoring of supplier
factories to ensure they comply with the law. But they acknowledge that ensuring compliance is
challenging in China.

A spokesman for the local Dongguan government here said that they were strictly enforcing the new
law. But in interviews, some factory owners acknowledged that they were seeking ways to get around
it, complaining that the law's regulations were too costly and cumbersome.

Lawyers say some local governments have issued their own competing rules or interpretations of the
law that weaken it, to aid factory owners.

``Many local governments want to develop their own versions of the law,'' says Liu Cheng, a
professor of law at Shanghai Normal University and one of the law's authors.

China's huge and complicated labor market has long thrived on cheap labor and lax regulation. In
recent years, labor rights advocates say they have seen incremental gains for workers. But they say
there are growing signs of labor abuse. They point to a string of recent cases, like one several
weeks ago in which police in southern China's Anhui province said they had freed 30 mentally
handicapped workers from what they called ``slave conditions'' in a brick kiln.

On the same day, police said a fire in the dormitory of an illegal factory in southern Guangdong
province killed 13 female workers and seriously injured four others.

A few weeks earlier, 7,000 workers went on strike at a factory that supplies some of the world's
biggest technology companies, saying they were being cheated on overtime wages and fed unsanitary
food.

Experts say cheating workers on wages, forcing them to log up to 200 hours of overtime a month and
denying them health benefits is commonplace in China.

Many factories are violating not just the new contract labor law, but also a 1994 law, which covered
a broader set of labor and wage practices, they said.

``The employment contract in many factories here is a mere scrap of paper,'' says Liu Kaiming,
director of the Institute of Contemporary Observation, a labor rights group in Shenzhen. ``Here is a
common trick: The factory signs contracts with 1,000 workers but actually they've hired 2,000. The
factory reports to the government saying they have 100 percent of their workers registered.''

Heather White, a consultant who has inspected factories in China for Tommy Hilfiger, Polo Ralph
Lauren and other big companies, says many exporters evade the law by subcontracting to so-called
shadow factories, which operate under illegal conditions.

``The market is penalizing anyone who complies with the law,'' she says, meaning their products are
more expensive. ``And so many companies are subcontracting'' to shadow factories.

Labor rights groups that specialize in sneaking into Chinese factories and documenting their flaws
say exporters' multinational clients are also responsible for their suppliers' practices.

``They are blatantly violating the labor law,'' says Charles Kernaghan, director of the National
Labor Committee, based in Pittsburgh, which last February issued a scathing report on a factory
making keyboards for big tech companies. ``They're forcing people to work 12-hour shifts. Their
overtime far exceeds the legal limit.''

But factory owners say that labor law enforcement has been weak and selective for years, and
changing the rules now could lead to chaos, drive up prices and force many factories out of
business.

``The government hasn't given us time to adjust,'' says Huang Zhenyuan, vice president of the Taiwan
Merchant Association of Dongguan, which represents thousands of factories. ``When we came to China
there was no legal environment. Now all has changed; it's too sudden.''

Because of the downturn, 20 million migrant workers have already lost their jobs, Beijing says. The
government recently put rules in place restricting factories from making large-scale layoffs without
giving the government notice.

But on an individual level, the struggle between having a job and economic security, and safety and
personal dignity can be wrenching.

Liu Pan, the worker crushed to death, was hired shortly after he had turned 15. He operated a giant
machine that turned out boxes in a plant that Disney concedes had recently passed third-party
audits. His salary was about \$175 a month.

Workers found his mutilated body stuck in the machine on the afternoon of April 5.

Michael Li, a senior manager at Yiuwah, says the accident was not a reflection of labor conditions
at the factory. He also said a Chinese government official helped manage the factory.

But China Labor Watch, a nonprofit group based in New York, says it investigated conditions at the
factory shortly after the death and found widespread violations of the labor law, including the
hiring of children as young as 13, forced overtime and the failure of many workers to sign labor
contacts.

In a statement, Disney said only about 5 to 15 percent of the goods produced at Yiuwah were made for
Disney and that Yiuwah had committed to correcting problems there.

``However, if improvement within acceptable and agreed upon time frames is not achieved,'' Disney
said it would stop doing business with the factory.

Yiuwah offered \$22,000 to compensate for Mr.~Liu's death, his family said.

Liu Hong, Mr.~Liu's father, does not even know how to begin to measure such compensation.

``I'm falling apart,'' he said as his wife tried to calm him. ``We are in the lowest class. So I
still don't know if it's the highest compensation. I still wonder, because a life, a young life, is
only worth \$22,000?''

He added, ``He was my only son, and he's the only grandson to my father.''

\section{Mind Your BlackBerry or Mind Your Manners}

For the first half-hour of the meeting, it was hardly surprising to see a potential client fiddling
with his iPhone, said Rowland Hobbs, the chief executive of a marketing firm in Manhattan.

At an hour, it seemed a bit much. And after an hour and a half, Mr.~Hobbs and his colleagues
wondered what the man could possibly be doing with his phone for the length of a summer blockbuster.

Someone peeked over his shoulder. ``He was playing a racing game,'' Mr.~Hobbs said. ``He did ask
questions, though, peering occasionally over his iPhone.''

But, Mr.~Hobbs added, ``we didn't say anything. We still wanted the business.''

As Web-enabled smartphones have become standard on the belts and in the totes of executives, people
in meetings are increasingly caving in to temptation to check e-mail, Facebook, Twitter, even
(shhh!) ESPN.com.

But a spirited debate about etiquette has broken out. Traditionalists say the use of BlackBerrys and
iPhones in meetings is as gauche as ordering out for pizza. Techno-evangelists insist that to ignore
real-time text messages in a need-it-yesterday world is to invite peril.

In Hollywood, both the Creative Artists Agency and United Talent Agency ban BlackBerry use at
meetings. Tom Golisano, a billionaire and power broker in New York State politics, said last week
that he pushed to remove Malcolm A.~Smith as the State Senate majority leader after the senator met
with him on budget matters in April and spent the time reading e-mail on his BlackBerry.

The phone use has become routine in the corporate and political worlds -- and grating to many. A
third of more than 5,300 workers polled in May by Yahoo HotJobs, a career research and job listings
Web site, said they frequently checked e-mail in meetings. Nearly 20 percent said they had been
castigated for poor manners regarding wireless devices.

Despite resistance, the etiquette debate seems to be tilting in the favor of smartphone use, many
executives said. Managing directors do it. Summer associates do it. It spans gender and generation,
private and public sectors.

A few years ago, only ``the investment banker types'' would use BlackBerrys in meetings, said Frank
Kneller, the chief executive of a company in Elk Grove Village, Ill., that makes water-treatment
systems. ``Now it's everybody.'' He said that if he spotted 6 of 10 colleagues tapping away, he knew
he had to speed up his presentation.

It is routine for Washington officials to bow heads silently around a conference table -- not
praying -- while others are speaking, said Philippe Reines, a senior adviser to Secretary of State
Hillary Rodham Clinton. Although BlackBerrys are banned in certain areas of the State Department
headquarters for security reasons, their use is epidemic where they are allowed.

``You'll have half the participants BlackBerrying each other as a submeeting, with a running
commentary on the primary meeting,'' Mr.~Reines said. ``BlackBerrys have become like cartoon thought
bubbles.''

Some professionals admitted that they occasionally sent mocking commentary about the proceedings,
but most insisted that they used smartphones for legitimate reasons: responding to deadline
requests, plumbing the Web for data to illuminate an issue under discussion or simply taking notes.

Still, the practice retains the potential to annoy. Joel I.~Klein, the New York City schools
chancellor, has gained such a reputation for checking his BlackBerry during public meetings that
some parents joke that they might as well send him an e-mail message. Few companies have formal
policies about smartphone use in meetings, according to Nancy Flynn, the executive director of the
ePolicy Institute, a consulting group in Columbus, Ohio. Ms.~Flynn tells clients to encourage
employees to turn off all devices.

``People mistakenly think that tapping is not as distracting as talking,'' she said. ``In fact, it
can be every bit as much if not more distracting. And it's pretty insulting to the speaker.''

Still, business can be won or lost, executives say, depending on how responsive you are to an e-mail
message. ``Clients assume they can get you anytime, anywhere,'' said David Brotherton, a media
consultant in Seattle. ``Consultants who aren't readily available 24/7 tend to languish.''

Playful electronic bantering can stimulate creativity in meetings, in the view of Josh Rabinowitz,
the director of music at Grey Group in New York, an advertising agency. In pitch meetings,
Mr.~Rabinowitz said, he often traded messages on his Palm Treo -- jokes, ideas, questions -- with
colleagues, ``things that you might not say out loud.''

The chatter tends to loosen the proceedings. ``It just seems to add to the productive energy,'' he
said.

But business relationships can be jeopardized. Lori Levine, the founder of Flying Television, a
talent-booking agency in Manhattan, said that in an effort to be environmentally sensitive she
instructed employees to take notes on BlackBerrys instead of paper during client meetings.

``Then I got a call from a client screaming that our vice president spent an hour on his BlackBerry
during a huge meeting,'' Ms.~Levine recalled. To soothe the client, Ms.~Levine read aloud the notes
the vice president had taken.

In Dallas, a college student sunk his chance to have an internship at a hedge fund last summer when
he pulled out a BlackBerry to look up a fact to help him make a point during his interview, then
lingered -- momentarily, but perceptibly -- to check a text message a friend had sent, said Trevor
Hanger, the head of equity trading at the hedge fund, who was helping conduct the interview.

Very few companies have policies on smartphone use in meetings, which leaves it up to employees to
feel their way across uncertain terrain.

To Jason Chan, a digital-strategy consultant in Manhattan, different rules apply for in-house
meetings (where checking BlackBerrys seems an expression of informal collegiality) and those with
clients, where the habit is likely to offend. There is safety in numbers, he added in an e-mail
message: ``The acceptability of checking devices is proportional to the number of people attending
the meeting. The more people there are, the less noticeable your typing will be.''

Beyond practical considerations, there is also the issue of image. In many professional circles,
where connections are power, making a show of reaching out to those connections even as co-workers
are presenting a spreadsheet presentation seems to have become a kind of workplace boast.

Mr.~Brotherton, the consultant, wrote in an e-mail message that it was customary now for
professionals to lay BlackBerrys or iPhones on a conference table before a meeting -- like
gunfighters placing their Colt revolvers on the card tables in a saloon. ``It's a not-so-subtle way
of signaling 'I'm connected. I'm busy. I'm important. And if this meeting doesn't hold my interest,
I've got 10 other things I can do instead.' ''

\section{Apple's Secrecy on Products and Top Executives}

SAN FRANCISCO -- Apple is one of the world's coolest companies. But there is one cool-company trend
it has rejected: chatting with the world through blogs and dropping tidbits of information about its
inner workings.

Few companies, indeed, are more secretive than Apple, or as punitive to those who dare violate the
company's rules on keeping tight control over information. Employees have been fired for leaking
news tidbits to outsiders, and the company has been known to spread disinformation about product
plans to its own workers.

``They make everyone super, super paranoid about security,'' said Mark Hamblin, who worked on the
touch-screen technology for the iPhone and left Apple last year. ``I have never seen anything else
like it at another company.''

But even by Apple's standards, its handling of news about the health of its chief executive and
co-founder, Steven P.~Jobs, who has battled pancreatic cancer and recently had a liver transplant
while on a leave of absence, is unparalleled.

Mr.~Jobs received the liver transplant about two months ago, according to people briefed on the
matter by current and former board members. Despite intense interest in Mr.~Jobs's condition among
the news media and investors, Apple representatives have declined to address the matter, reciting
with maddening discipline only that Mr.~Jobs is due back at the company by the end of June.

Mr.~Jobs was actually at work on Apple's sprawling corporate campus on Monday, according to a person
who saw him there. Company representatives would not say whether he had returned permanently.

Even senior officials at Apple fear crossing Mr.~Jobs. One official, who is normally more open, when
asked for a deep-background briefing about Mr.~Jobs's health after the news of the transplant had
become public, replied: ``Just can't do it. Too sensitive.''

Secrecy at Apple is not just the prevailing communications strategy; it is baked into the corporate
culture. Employees working on top-secret projects must pass through a maze of security doors,
swiping their badges again and again and finally entering a numeric code to reach their offices,
according to one former employee who worked in such areas.

Work spaces are typically monitored by security cameras, this employee said. Some Apple workers in
the most critical product-testing rooms must cover up devices with black cloaks when they are
working on them, and turn on a red warning light when devices are unmasked so that everyone knows to
be extra-careful, he said.

Apple employees are often just as surprised about new products as everyone else.

``I was at the iPod launch,'' said Edward Eigerman, who spent four years as a systems engineer at
Apple and now runs his own technology consulting firm. ``No one that I worked with saw that
coming.''

Mr.~Eigerman was fired from Apple in 2005 when he was implicated in an incident in which a co-worker
leaked a preview of some new software to a business customer as a favor. He said Apple routinely
tries to find and fire leakers.

Philip Schiller, Apple's senior vice president for marketing, has held internal meetings about new
products and provided incorrect information about a product's price or features, according to a
former employee who signed an agreement not to discuss internal matters. Apple then tries to track
down the source of news reports that include the incorrect details.

Five years ago, Apple took its obsession with secrecy to the courts. It sued several bloggers who
had covered the company, arguing that they had violated trade-secret laws and were not entitled to
First Amendment protections. A California appeals court ruled for the bloggers, and the company had
to pay \$700,000 in legal fees.

Apple also sued a blog called Think Secret and settled the case for an undisclosed amount, but as
part of the settlement that blog shut down.

Regis McKenna, a well-known Silicon Valley marketing veteran who advised Apple on its media strategy
in its early days, said the culture of secrecy had its origin in the release of the first Macintosh,
which competitors like Microsoft and Sony knew about before it was unveiled.

``It really started around trying to keep the surprise aspect to product launches, which can have a
lot of power,'' Mr.~McKenna said.

He added: ``But what most people don't understand is that Steve has always been very personal about
his life. He has always kept things close to the vest since I've known him, and only confided in
relatively few people.''

Apple's decision to severely limit communication with the news media, shareholders and the public is
at odds with the approach taken by many other companies, which are embracing online outlets like
blogs and Twitter and generally trying to be more open with shareholders and more responsive to
customers.

``They don't communicate. It's a total black box,'' said Gene Munster, an analyst at Piper Jaffray
who has covered Apple for the last five years.

Mr.~Munster said he jokes with other colleagues covering the company about how Apple routinely
``jams the frequencies,'' or gives them misinformation to throw them off the scent of a new product
or other news it hopes to keep confidential. Four years ago, he said, a senior Apple executive
directly told him the company had no interest in developing a cheap iPod with no screen. Soon after,
the company released just that: the iPod Shuffle.

For corporate governance experts, and perhaps federal regulators, the biggest question is whether
Mr.~Jobs's approach has led to violating laws that cover what companies must disclose to the public
about the well-being of their chief executive.

On that key issue, the experts are divided. Some believe Apple did not need to disclose Mr.~Jobs's
liver transplant because Mr.~Jobs was on a leave of absence and had passed responsibility for the
day-to-day operations of the company to the chief operating officer, Timothy Cook.

Other governance experts argue that the liver transplant now makes one of Apple's assertions from
January -- that Mr.~Jobs was suffering only from a hormonal imbalance -- seem like a deliberate
mistruth, unless Mr.~Jobs's health condition suddenly deteriorated. Of course, no one knows enough
to say definitively.

Most governance experts do seem to agree on one point: that the secrecy that adds surprise and
excitement to Apple product announcements is not serving the company well in other areas.

``In this environment, where transparency is critical, the more information you give the marketplace
the better,'' said Charles Elson, director of the John L.~Weinberg Center for Corporate Governance
at the University of Delaware. ``For a technology company that views itself as innovative, it's a
little odd that they are getting a reputation for lack of transparency.''

Apple's stock dropped \$2.11 to \$137.37 on Monday amid a larger market sell-off. And the company
did, in fact, have something to reveal: it said it had sold a million units of its new iPhone 3G S
over the weekend, well above analysts' forecasts.

\section{Nortel Falls To Telecom's Tectonic Shift}

Earlier this month, when I wrote about Telecom's Titanic Shifts and the decline of the once mighty
service providers, in passing I noted the slow-mo descent of Western equipment makers. With the
mega-growth registered by non-Western carriers as dominant equipment buyers, we have seen the rise
of two hardware giants: Huawei and ZTE.

Their rising fortunes have been mirrored by sharp declines for hardware makers based in Europe and
North America. That point was brought home over the weekend when Nortel announced that it was
selling CDMA and LTE pieces of its wireless equipment business to Nokia Siemens Networks (NSN) for
about \$650 million. It was a tactical admission by the once-mighty Nortel that it had little or no
room in tomorrow's telecom world.

I have followed Nortel for a long time. The company rose and fell with the telecom bubble. It got
Cisco envy -- it bought companies for outrageous prices only to see the deals go nowhere. The hubris
of its managers turned this once proud house of innovation into a footnote. The company started to
chase new markets too late and relied too much on North American market for too long. A lot of that
can be blamed on lack of management.

How long before we see other big names in telecom take a dirt nap? Will Alcatel-Lucent be next? In a
note to its clients, UBS Research points out that Alcatel-Lucent is going to face considerable
pressure in North America, as NSN will indeed be super aggressive. Given how aggressive Huawei has
been in its bid to gain market share, the jury is still out even on NSN's long-term sustainability.

\section{Apple Retail Store Success: It Ain't Rocket Science}

In May of 2001, Apple opened its first retail outlet store in Tysons Corner, Va. It was immediately
met by critics and the typical roster of Apple-haters with predictions of doom and massive financial
losses. Eight years later, Apple has more than 255 retail stores worldwide, and are the darling of
the retail computer industry. So what makes them so successful?

While the Mac, iPod, and iPhone are what makes Apple successful, the formula for the retail stores'
success isn't so obvious until you actually visit an Apple store. Sure, great products sell
themselves, but Apple has made sure that its products have the best opportunity to do so by creating
an environment that offers virtually no reason for a buyer to go elsewhere.

The Typical Computer-Buying Experience

When I think of how I used to purchase computers in the past at stores like CompUSA, Circuit City,
Computer City (all three now out of business), I do so with little in the way of fond memories.
Dimly-lit stores with out-of-date computers, software, and hardware litter my thoughts. CompUSA in
particular was famous for the horrible way in which it showcased their hardware offerings. Most
computers were dirty, had broken keyboards (if a keyboard was attached to the computer at all), no
monitor attached, and usually had no software installed -- which mattered little, because if you
actually found a computer that worked, it was typically password protected and there was rarely a
sales rep to be found.

Even today, most stores like Staples, Best Buy, OfficeMax, and Fry's Electronics make the computer
buying experience as difficult as possible. Desktop computers typically sit idle on shelves high out
of reach and completely lifeless. Laptops are barred, bolted, strapped and nailed down to a low
display shelf, with as many of them as the store can squeeze in. One person looking at a laptop in
the middle of the isle effectively prevents anyone from looking at that laptop, and the one on each
side of it. Price tags on the shelf rarely coincide with the computer sitting right above it, and
offer little in the way of computer specs.

You can't pick them up or comfortably test the keyboard in a natural position, and just like the
previously mentioned stores, they either have no power running to them, or are password protected to
prevent you from actually using them. The few working models you do find have little-to-no software
installed, and none have Internet access.

The Apple Experience

It's really not difficult to understand why Apple chose to open its own stores. Steve Jobs wanted a
better environment for buyers to experience the Macintosh. And Apple pulled it off masterfully. It's
all about the experience.

Visitors at the Apple Store at the Scottsdale Quarter

Spacious interior – The isles are wide, allowing patrons to move freely about the store. In fact,
Apple almost encourages loitering with the amount of empty space and abundance of stools and tables.
Plenty of hardware – Apple puts as many units of each model of computer it offers on display as they
can comfortably fit. Other than launch days and special events, you'll likely have no trouble
getting at the hardware you want to check out. The Scottsdale Quarter store I recently visited had
no less than 20 iPod touch and iPhone's each on display, at least 20 laptops, and nearly a dozen or
more iMacs and MacPros. In working order – Every single device Apple sells in their stores is
working. And they don't limit the access to the devices with passwords that require you to find an
available sales associate. Internet Access – All computers, iPod touches, and iPhones are connected
to the store's Wi-Fi network. You can get on the web and check out your favorite sites on any
device. This is such an obvious sales tool that most other stores simply ignore. Easily accessible
display units at the Apple store

Accessible – You buy a laptop, iPod, or iPhone to take on the road. Comfort, weight and the overall
feel of the device is important. Apple knows this, and makes it easy to handle the display models.
Unlike most other stores, they're not bolted down. Prices are clearly marked. Realistic Expectations
– Every Mac in the store has the latest OS installed, as well as iLife and iWork. You can get a real
feel for how the apps you typically use will perform on any particular machine. In many cases, Apple
even installs Final Cut, and other third-party software such as Photoshop, Quark XPress, and
Microsoft Office on some Macs. This is especially useful for high-end users looking to check out the
latest Mac Pro. Beyond the Macs

Along with their own hardware and software, Apple carefully chooses which third-party hardware and
software to sell in its stores. You won't find every obscure make and model of input devices,
cameras, printers, speakers and software. Apple goes out of its way to select the best models that
work 100 percent out of the box with the Mac. And like its own products, Apple puts them on display
so you can easily try them out to see how they work. By limiting the amount of software and hardware
it sells from other vendors, it ensures a pleasant and positive experience for the buyer.

Then of course there's the check-out process and the Genius Bar. Most computer stores have four or
more check-out registers, but only have one person working them. At the Apple store, you simply look
for an available associate who can run your credit card through a handheld device and have your
receipt emailed to you. Simple! And if you have any questions about your purchase, or have issues
with an Apple product you already own, the Genius Bar offers free advice and warranty work -- take
that, GeekSquad!

Though the Apple retail stores probably aren't as close to your home or office as the nearest
Staples, it's certainly worth the drive if you have one near you. The entire experience is much more
satisfying than any other retail experience I've come across. You don't have to be a rocket scientist to see that Apple simply makes it a place you want to shop at, rather than one you have to.

\section{Obama Assails Iron Fist in Iran Against Protesters}

WASHINGTON -- President Obama hardened his tone toward Iran on Tuesday, condemning the government
for its crackdown against election protesters and accusing Iran's leaders of fabricating charges
against the United States.

In his strongest comments since the crisis erupted 10 days ago, Mr.~Obama used unambiguous language
to assail the Iranian government during a news conference at the White House, calling himself
``appalled and outraged by the threats, beatings and imprisonments of the past few days.''

He praised what he called the courage and dignity of the demonstrators, especially the women who
have been marching, and said that he had watched the ``heartbreaking'' video of a 26-year-old
Iranian woman whose last seconds of life were captured by video camera after she was shot on a
Tehran street.

``While this loss is raw and extraordinarily painful,'' he said, ``we also know this: Those who
stand up for justice are always on the right side of history.''

Yet beyond muscular words, Mr.~Obama has limited tools for bringing pressure to bear on the Iranian
government, which for years has been brushing off international calls for it to curb its nuclear
program.

After the news conference, administration officials said there was little they could do to influence
the outcome of the confrontation between the government and the protesters. And more so now than
even a few days ago, they said, the prospects for any dialogue with Iran over its nuclear program
appear all but dead for the immediate future, though they held out hope that Iran, assuming it has a
stable government, could respond to Mr.~Obama's overtures later in the year.

At home, Mr.~Obama has been under intense pressure, especially from conservatives, to align the
United States more forcefully with the protesters. On Tuesday, he dismissed suggestions that he had
changed his tone toward Iran in response to critical comments from Senator John McCain of Arizona
and other Republicans.

In sometimes testy exchanges with reporters at the news conference, Mr.~Obama defended himself,
contending that even the moderate tone he had struck previously had been twisted by Iran's
government to suggest that the protests had been engineered by the United States.

``They've got some of the comments that I've made being mistranslated in Iran, suggesting that I'm
telling rioters to go out and riot some more,'' Mr.~Obama said, referring to accounts that the White
House said surfaced late last week and over the weekend. ``There are reports suggesting that the
C.I.A. is behind all this. All of which is patently false. But it gives you a sense of the narrative
that the Iranian government would love to play into.''

But after the crackdown over the weekend that left an untold number of protesters dead -- and after
the wide dissemination of the video of the last moments of Neda Agha-Soltan, the Iranian woman who
appeared to lock eyes with the camera as she died after being shot -- White House officials decided
that Mr.~Obama had to take a tougher stand.

``The situation looked very different on Saturday than it did when he first spoke in the Oval Office
a week ago,'' one of Mr.~Obama's media advisers said.

``It was the bloodshed'' that led to the change in tone, he said.

While Mr.~Obama did not rule out the possibility of engaging with Iran over the nuclear issue,
administration officials and European diplomats say that the door to talks has all but closed, at
least for now.

``I think that under these circumstances, no one is going to be able to pursue anything because
there is nothing to pursue,'' said Trita Parsi, the president of the National Iranian American
Council, who has been consulting with White House officials ``on a daily basis,'' he said, about the
unfolding situation in Iran.

Mr.~Parsi said that all past assumptions about where Iran was headed had been cast aside by the
disputed election results and the response of the protesters.

Administration officials acknowledged that after reading reams of intelligence reports, watching
videos of the street demonstrations and absorbing the trickle of intelligence from Iran, they were
unable to predict how the protests might turn out.

During the news conference, Mr.~Obama maintained that he had been consistent in his tone toward Iran
all along. ``As soon as violence broke out -- in fact, in anticipation of potential violence -- we
were very clear in saying that violence was unacceptable, that that was not how governments operate
with respect to their people,'' Mr.~Obama said.

But the language Mr.~Obama used on Tuesday was more forceful and less ambiguous than his previous
statements. In an interview with CNBC and The New York Times last week, he said that as far as
America's national interests were concerned, there was not much difference between Iran's president,
Mahmoud Ahmadinejad, and his challenger, Mir Hussein Moussavi.

In his first public comment on the situation on June 15, Mr.~Obama said he was troubled by the
postelection violence and called on Iran's leaders to respect free speech and the democratic
process. The next day, on June 16, he said he had ``deep concerns'' about the elections but also
said that it would be counterproductive for the United States ``to be seen as meddling.''

In the internal discussions at the White House about how to handle Iran, Mr.~Obama's aides are
clearly struggling with how to reconcile two different goals: supporting a nascent, unpredicted
movement in the streets that could weaken the country's top clerics, and following the diplomatic
mixture of pressure and diplomacy that Mr.~Obama settled on months ago as a strategy to halt Iran's
nuclear work.

The protests, administration officials said, create the first possibility in 30 years that the
mullahs' grip on Iran might be loosened. Even if the street protests are put down, one official
said, ``a fissure has opened up that cannot be completely closed.''

Clearly those events took the administration by surprise: none of the possibilities for the election
that were laid out for Mr.~Obama a month ago, one official said, included the possibility of a
violently disputed election.

Yet in the long run, Mr.~Obama's aides say, they are not certain that the protests will change the
fundamental calculus about the risks Iran poses to its neighbors or the United States. One of
Mr.~Obama's strategists noted that ``one has to be concerned that while all this is happening the
centrifuges are still spinning,'' a reference to the machines used to enrich uranium.

\end{document}
